\chapter{Contenidos y secciones principales}
Dependiendo del tipo de informe escrito y la naturaleza de la información, la estructura del documento puede variar. Esta sección presenta los elementos mínimos que debe contener un informe escrito. Además, se dan ejemplos de los contenidos que deberían tener TFT, TFG, tesis de magíster y tesis de doctorado. 

\section{Contenidos mínimos recomendados}
\begin{enumerate}
    \item Introducción: Ofrece una visión global del tema tratado, destacando su importancia y pertinencia dentro del campo de estudio. Establece el contexto y motiva al lector a profundizar en el documento. 
    \item Descripción del problema: Detalla el problema específico que se propone abordar en el documento. Explica de manera clara y concisa, señalando por qué es un problema relevante y digno de investigación o estudio.
    \item Objetivos: Define los objetivos generales y específicos del trabajo. Los objetivos generales señalan la meta amplia del estudio, mientras que los específicos detallan los pasos concretos para alcanzar dicha meta, marcando el camino de la investigación o desarrollo. 
    \item Marco teórico/referencial: Repasa las teorías, modelos, y estudios previos relacionados con el tema de desarrollo o investigación. Sirve para contextualizar el trabajo dentro del conocimiento existente y establecer el fundamento teórico sobre el cual se construye la propuesta. 
    \item Metodología y plan de trabajo: En esta parte, se describe detalladamente cómo se lleva a cabo la investigación. Incluye la elección de métodos de investigación, diseño de experimentos, técnicas de recopilación y análisis de datos, y herramientas utilizadas. Esto puede abarcar desde la definición de algoritmos hasta el uso de frameworks específicos para el análisis de datos. 
    \item Propuesta de desarrollo: Describe de forma detallada la solución o enfoque propuesto para abordar el problema identificado. Incluye metodologías, técnicas, herramientas, y cualquier otro recurso considerado en el desarrollo del trabajo.
    \item Validación y pruebas (si corresponde): Explica cómo se ha llevado a cabo la validación y prueba de la solución propuesta. Dependiendo del tipo de trabajo, esto puede incluir experimentos, análisis de datos, estudios de caso, entre otros, para demostrar la eficacia y viabilidad de la propuesta. 
    \item Conclusiones: Presenta una síntesis de los principales hallazgos y resultados obtenidos, las implicaciones de estos dentro del campo de estudio, y las recomendaciones para investigaciones o desarrollos futuros. Esta sección refleja el aporte del trabajo al conocimiento existente y su relevancia práctica o teórica. 
\end{enumerate}